\chapter*{Abstract}
\addcontentsline{toc}{chapter}{Abstract}
\markboth{Abstract}{Abstract}

Large language models have shown strong potential for instruction understanding and high-level reasoning in human-robot interaction. Their integration into a robot nevertheless requires more than a planner call: dialogue state, symbolic grounding, capability descriptions, execution admission, robot-facing skills, feedback, and user-facing reporting must remain coherent across a distributed runtime. This thesis presents the design, integration, and validation of a modular language-model-enabled interaction and execution stack for the NAO robot.

The system separates dialogue management, semantic routing, symbolic grounding, task planning, deterministic orchestration, and skill execution. A chatbot interprets each turn and selects a conversational, knowledge, or execution route. A planner generates structured plans over a reviewed skill registry, while an orchestrator admits requests, validates plans, dispatches real or simulated skills, and returns typed execution feedback. A knowledge and perception path provides compact grounded context without allowing planning-time state to stand in for execution evidence.

The contribution is an inspectable NAO interaction stack and a contract-based validation method for its principal boundaries. The validation follows complete user turns as well as isolated planner and deterministic fake-skill profiles. It examines route correctness, plan validity, entity and role preservation, execution lineage, symbolic postconditions, failure handling, and exactly-once user-facing closure. Direct atomic execution is retained as an orchestrator control rather than treated as a causal baseline for multi-step planning.

\vspace{10pt}
\textbf{Keywords:} Human-Robot Interaction, ROS 2, Large Language Models, Task Planning, NAO Robot, Symbolic Grounding, Deterministic Execution, Runtime Validation.
